% Chapter 7
\chapter{Conclusion and Future Work}
\label{Chapter7}
\lhead{Chapter 7. \emph{Conclusion and Future Work}}

\section{Conclusion}
This thesis presented Magnavis, a software framework for magnetic observatory data that links ingestion, recurrent forecasting, residual-based anomaly detection, and optional perturbation-direction visualisation. The main scientific contribution of the current work is a controlled benchmark under zero-historic, predict-only operation, comparing statistical smoothers, a session-trained attention bidirectional LSTM, and pretrained GRU and LSTM models on short, long, and synthetic ground-truth layouts.

Three conclusions follow from the experiments. First, detector multiplier $k$ must be chosen together with the predictor: on long-duration ground truth, offline smoothers fail at moderate $k$, whereas neural predictors can remain useful. Second, pretrained LSTM delivers roughly flat 85\% recall on Apr-27 for $k=2$--$5$; pretrained GRU matches LSTM at $k=3$ but falls to about 63\% for $k \ge 4$ because interior residuals shrink when the model tracks a sustained offset. Closed-loop autoregressive rolling does not remove that gap. Third, for operational use on long campaigns, LSTM at $k \approx 3$ is the most robust setting identified here; GRU at high $k$ on long ground truth requires retraining or careful validation.

The direction-finding module provides a geometric characterisation of tri-axial perturbations and optional two-station triangulation. Field experience shows that reported bearings may disagree with known source positions unless baselines, alignment, and timing are carefully controlled.

\section{Limitations}
\begin{itemize}
    \item Point-level metrics on fixed CSV exports may not capture event-level detection delay or operator workflow in live mode.
    \item Pretrained weights reflect a specific training archive and sensor; transfer to other sites is not guaranteed.
    \item Direction finding reports the perturbation vector, not a unique source location.
    \item Database credentials and session logs in the prototype should be externalised before production deployment.
\end{itemize}

\section{Future Work}
\begin{enumerate}[label=(\alph*)]
    \item Event-level recall and detection latency on manual anomaly intervals.
    \item Systematic GRU retraining studies (absolute versus delta targets) and checkpoint selection criteria.
    \item Calibration of direction estimates using controlled source runs at known bearings.
    \item Dipole or multi-point inversion using vector anomalies at two observatories.
    \item Extension of the benchmark to additional OBS sensors and months of archive data.
    \item Integration of the evaluation harness into continuous integration for regression testing after model updates.
\end{enumerate}

\section{Closing Remarks}
Magnavis demonstrates that a single open pipeline can support both interactive monitoring and rigorous offline comparison of forecasting and detection choices. The repository, bundled models, and benchmark summaries provide a reproducible basis for thesis examination and for further research on static observatory magnetic anomaly detection at IIT Kanpur.
